\chapter{Theoretical Background}
\label{ch:background}

\section{Preference Optimization: The Bradley--Terry Model}
\label{sec:bg-bradley-terry}

The Bradley--Terry model \citep{bradley1952rank} is a standard probabilistic model
of pairwise comparison.
It assumes that every item $i$ has an underlying, unobserved scalar quality $r_i$,
and that the probability of item $i$ being preferred over item $j$ in a pairwise
comparison is
\begin{equation}
  P(i \succ j) = \sigma(r_i - r_j),
  \qquad
  \sigma(x) = \frac{1}{1 + e^{-x}}.
  \label{eq:bradley-terry}
\end{equation}
Given a dataset of observed pairwise preferences $(w, l)$, where $w$ (the
``winner'') was judged preferable to $l$ (the ``loser''), the model parameters are
fit by maximizing the likelihood of the observed comparisons under
Equation~\ref{eq:bradley-terry}. Here the parameters belong to a neural network
$R_\theta$ that maps an (article, summary) pair to a scalar reward, and maximizing
the likelihood is equivalent to minimizing the negative log-likelihood
\begin{equation}
  \mathcal{L}(\theta) = - \, \mathbb{E}_{(w, l) \sim \mathcal{D}}
  \Big[ \log \sigma\big( R_\theta(w) - R_\theta(l) \big) \Big],
  \label{eq:bt-loss}
\end{equation}
where $\mathcal{D}$ is the preference dataset, holistic or GCA in this thesis.
Equation~\ref{eq:bt-loss} is exactly the loss implemented in \url{src/reward_model/train.py}
(\url{bradley_terry_loss}), and it belongs to the same family of
reward-modeling loss used inside RLHF pipelines more broadly
\citep{stiennon2020learning, ouyang2022training}. The difference here is that the
reward model $R_\theta$ is not an intermediate component on the way to fine-tuning a
generation policy; it is the final object of comparison. The research question this
thesis asks, whether GCA produces a more learnable preference signal than holistic
scoring, is operationalized directly: does a Bradley--Terry reward model trained on
GCA-derived preferences achieve higher held-out pairwise accuracy than one trained
on holistic-derived preferences, under an otherwise identical training recipe?

This corresponds to the inverse reinforcement learning perspective: given a set
of preference demonstrations, the task is to recover the latent reward function
that best explains them. Comparing two reward models trained on two differently
constructed preference datasets is then a controlled test of which construction
procedure yields the more consistently recoverable latent signal, without
requiring a downstream generation policy to be fine-tuned and re-evaluated
(Section~\ref{sec:bg-dpo}).

The architecture used throughout this thesis (\texttt{BradleyTerryRewardModel} in
\url{src/reward_model/train.py}) encodes the concatenation of a truncated source
article and a candidate summary with a pretrained transformer encoder, mean-pools
the encoder's final hidden states over non-padding tokens, and maps the pooled
representation to a scalar reward through a single linear layer. Reward-model
accuracy is defined as the fraction of held-out pairs for which the model correctly
ranks the chosen summary above the rejected one, i.e.\ $R_\theta(w) > R_\theta(l)$
(\texttt{\_evaluate} in \url{src/reward_model/train.py}); this is the metric
reported throughout Chapter~\ref{ch:results}.

\section{Direct Preference Optimization (DPO)}
\label{sec:bg-dpo}

Direct Preference Optimization (DPO) replaces RLHF's reward-model-plus-reinforcement-learning
pipeline with a single closed-form supervised objective on a policy, defined
relative to a fixed reference policy \citep{rafailov2023dpo}. It is included
here only as background on the study's earlier design: DPO was the objective
originally proposed for this thesis and produced the early results reported in
Appendix~\ref{app:exploratory}, Section~\ref{sec:results-early}, but it is not
used in the final experiments (Chapter~\ref{ch:methodology},
Section~\ref{sec:methodology-evolution}).

The relevant point is structural. The DPO objective consumes exactly the same
pairwise preference pairs used to train the reward models in
Section~\ref{sec:bg-bradley-terry}, but converts them directly into a policy
update, so the quality of the preference data and the behavior of the optimizer
both influence the resulting summaries. Separating those two influences is what
motivated measuring the preference data on its own, which is what the final
study does.

\section{AlignScore: Factual Consistency Scoring}
\label{sec:bg-alignscore}

AlignScore is a unified alignment function trained across a mixture of natural
language processing tasks recast into a single ``information alignment'' format
\citep{zha2023alignscore}. Given a context $c$, the source article in this thesis,
and a claim $a$, a candidate summary or a single sentence from one, AlignScore
produces a scalar score in $[0, 1]$ estimating the degree to which $a$ is supported
by $c$. This thesis uses AlignScore unmodified as the sole factuality judge for
both preference-construction regimes
(\url{src/judging/reward_model_judge.py}, class \texttt{RewardModelJudge}).

AlignScore exposes several evaluation modes, of which four are used in this thesis
(the \texttt{alignscore-evaluation-mode} argument of
\url{build_reward_preferences.py}):

\begin{itemize}
  \item \texttt{nli\_sp}: the context is first split into sentences before scoring,
    and the claim is scored against this sentence-split representation via
    AlignScore's NLI (three-way entailment/neutral/contradiction) output head. This
    was the default mode inherited from the earliest AlignScore-based experiments
    in the project.
  \item \texttt{nli}: the same NLI output head, without the sentence-splitting
    preprocessing step.
  \item \texttt{bin\_sp} / \texttt{bin}: the same two preprocessing variants, but
    using AlignScore's binary (aligned / not aligned) output head instead of the
    three-way NLI head.
\end{itemize}

Chapter~\ref{ch:results}, Section~\ref{sec:results-mode-sweep} reports an empirical
sweep across all four modes. The choice of mode is associated with substantial
changes in the outcome, alongside the GCA aggregation formula
(Chapter~\ref{ch:methodology}, Section~\ref{sec:bg-gca}). The \texttt{nli} mode produced the largest GCA
advantage in that sweep and was selected as the intended final
configuration for the validated $n=1{,}000$ campaign and the
scale-robustness checks (Appendix~\ref{app:exploratory},
Section~\ref{sec:results-final-campaign}), though
Chapter~\ref{ch:implementation}, Section~\ref{sec:impl-preferences} reports
that the confirmatory twenty-run campaigns were, due to an implementation
oversight, actually trained under \texttt{nli\_sp} instead. This should be read as a
judge-\emph{configuration} finding specific to this pipeline and this task, not as
a general property of sentence-level aggregation independent of the underlying
judge; Chapter~\ref{ch:discussion} discusses what can and cannot be concluded from
it.

\section{Statistical Validation: Bootstrap Confidence Intervals and the Wilcoxon Signed-Rank Test}
\label{sec:bg-stats}

Two statistical tools are used throughout Chapter~\ref{ch:results} to assess
whether an observed GCA-versus-holistic accuracy gap is distinguishable from
noise.

\textbf{Bootstrap confidence intervals.} The bootstrap \citep{efron1979bootstrap}
estimates the sampling distribution of a statistic by resampling the observed data.
Given a set of paired fold-level accuracy
differences $\{d_1, \ldots, d_m\}$, where $d_k = \mathrm{acc}_{\text{GCA}, k} -
\mathrm{acc}_{\text{Holistic}, k}$ for cross-validation fold $k$, pooled across
independent training seeds, a bootstrap 95\% confidence interval for the mean
difference $\bar{d}$ is constructed by resampling $\{d_1, \ldots, d_m\}$ with
replacement $B$ times, computing the mean of each resample, and taking the 2.5th
and 97.5th percentiles of the resulting distribution of means. This thesis uses $B
= 10{,}000$ resamples with a fixed NumPy random seed of 42.

\textbf{Wilcoxon signed-rank test.} As a complementary non-parametric test suited
to a small number of paired, non-normally-distributed differences, the two-sided
Wilcoxon signed-rank test \citep{wilcoxon1945individual} is reported on the same paired differences. Unlike a
paired $t$-test it does not assume normality, which is appropriate given the small
number of observations and the fact that fold accuracies are bounded proportions.

\textbf{An important qualification: the unit of replication.} Both procedures
assume that the observations they are applied to are independent, and this
determines which level the tests may legitimately be applied at. Fold-level
differences are \emph{not} independent: the five folds within a run partition a
single dataset and therefore share most of their training data, so applying
either procedure to 30 such observations treats correlated measurements as
independent replications, which understates uncertainty and produces intervals
that are too narrow and $p$-values that are too small. Run-level differences,
by contrast, come from separately seeded training runs and satisfy the
assumption.

This thesis therefore reports two kinds of figure, and the distinction should be
kept in view when reading Chapter~\ref{ch:results}. Fold-level intervals and
$p$-values, which appear only for the superseded six-run campaign
(Appendix~\ref{app:exploratory}, Section~\ref{sec:results-final-campaign}), are
anti-conservative and are presented as exploratory summaries of spread, never as
hypothesis tests. All confirmatory statistics (the twenty-run campaigns at each
dataset size, their bootstrap intervals, Wilcoxon tests, effect sizes, and
equivalence tests) are computed over run-level differences, where the
independence assumption holds. Chapter~\ref{ch:discussion},
Section~\ref{sec:disc-stats-caveat} sets out what each level can support.

\textbf{Effect size.} A $p$-value indicates whether a difference is
distinguishable from zero, not how large it is, so both are reported. For paired
differences $\{d_1, \ldots, d_n\}$ the standardized effect size is Cohen's $d =
\bar{d} / s_d$, the mean difference expressed in units of its own standard
deviation \citep{cohen1988statistical}. Alongside it, the matched-pairs
rank-biserial correlation is reported as the effect size native to the Wilcoxon
test: it is the difference between the sums of positive and negative signed
ranks, divided by their total, and therefore ranges from $-1$ (every run favors
one condition) to $+1$ (every run favors the other), capturing consistency of
direction independently of magnitude.

\textbf{Equivalence testing.} A non-significant result does not by itself
establish that an effect is absent; it may equally reflect insufficient data.
Establishing absence requires reversing the burden of proof, which is what the
two one-sided tests (TOST) procedure does \citep{lakens2017equivalence}. Given an
equivalence bound $\pm\delta$ representing the smallest effect that would be
practically meaningful, TOST tests the two null hypotheses that the true effect
is at least $+\delta$ and at most $-\delta$. Rejecting both places the true
effect inside $(-\delta, +\delta)$ with the stated confidence, which is positive
evidence of a negligible effect rather than mere failure to detect one. This
procedure is applied in Chapter~\ref{ch:results},
Section~\ref{sec:results-effect-sizes} to the $n=5{,}000$ and $n=10{,}000$
campaigns, where the substantive claim is precisely that no advantage remains.

Pooling across runs rather than reporting a single run is nonetheless deliberate,
and responds to an early observation (Appendix~\ref{app:exploratory},
Section~\ref{sec:results-alpha0-validation}): a promising single-run result, a
$+1.4$ percentage-point advantage from one hyperparameter configuration, did not
reproduce when the same configuration was run again. Single runs at this dataset
scale are therefore not reliable evidence on their own, whatever statistics are
computed from their folds.
